\documentclass[12pt]{article}
\usepackage[T1]{fontenc}
\usepackage[utf8]{inputenc}
\usepackage{lmodern}
\usepackage{textcomp}
\usepackage{enumitem}
\usepackage{geometry}
\geometry{margin=1in}
\begin{document}
\begin{center}
Answer Key - History - K-12 Level Assessment 12 \
\small (Answers are ordered exactly as in the question paper.)
\end{center}

\section*{Multiple Choice}
\begin{enumerate}[label=\arabic*.]
\item \textbf{Answer:} C
\\ \textit{Options:} A) legumes, sarsens; B) lintels, stelae; C) sarsens, lintels; D) trilithon, tholoi
\\ \textit{Note:} The correct answer is "sarsens" for the upright stones and "lintels" for the horizontal stones because sarsens are the large sandstone blocks that make up the vertical elements of Stonehenge, while lintels are the horizontal stones that connect them at the top.
\item \textbf{Answer:} C
\\ \textit{Options:} A) deduction.; B) theory.; C) hypothesis.; D) scientific method.
\\ \textit{Note:} A hypothesis is a proposed explanation for an observed phenomenon that is formulated based on inductive reasoning, which involves generalizing from specific empirical observations.
\item \textbf{Answer:} B
\\ \textit{Options:} A) Italians.; B) Asians.; C) Australians.; D) Melanesians.
\\ \textit{Note:} Giovanni de Verrazzano described the native people of Rhode Island as resembling Asians due to their physical features, which he observed during his exploration of the region in the 16 th century.
\item \textbf{Answer:} C
\\ \textit{Options:} A) the feet; B) the skull; C) the foramen magnum; D) the ulna
\\ \textit{Note:} The foramen magnum is the opening at the base of the skull where the spinal cord connects to the brain, and its position determines the alignment of the head and spine, which influences whether a primate walks on all fours (quadruped) or stands upright on two legs (biped).
\item \textbf{Answer:} A
\\ \textit{Options:} A) Hang-t'u; B) Scapulimancy; C) Lung-shan; D) Qin manufacture
\\ \textit{Note:} Hang-t'u refers specifically to the ancient Chinese technique of using stamped or pounded earth for constructing durable house walls and defensive structures, making it the correct answer.
\end{enumerate}

\section*{True / False}
\begin{enumerate}[label=\arabic*.]
\setcounter{enumi}{5}
\item \textbf{Answer:} True
\\ \textit{Options:} A) True; B) False
\\ \textit{Note:} The statement is true because the ratio of carbon isotopes, specifically 13 C, in bones can indicate the type of plants consumed, as different plants have distinct carbon isotopic signatures related to their photosynthetic pathways.
\item \textbf{Answer:} False
\\ \textit{Options:} A) True; B) False
\\ \textit{Note:} The sagittal crest is actually a bony ridge located on the top of the skull, not above the upper jaw, and it serves as an attachment point for strong jaw muscles in some species.
\item \textbf{Answer:} False
\\ \textit{Options:} A) True; B) False
\\ \textit{Note:} The statement is false because Neandertals actually had larger noses compared to modern Homo sapiens, which were likely adaptations for cold climates, while their brain sizes were comparable, but not necessarily larger, than those of contemporary humans.
\item \textbf{Answer:} False
\\ \textit{Options:} A) True; B) False
\\ \textit{Note:} The Kwakiutl culture was primarily located along the Pacific Northwest coast of North America, particularly in present-day British Columbia, rather than the Ohio River valley.
\item \textbf{Answer:} True
\\ \textit{Options:} A) True; B) False
\\ \textit{Note:} Diffusionism is a theoretical framework in history that posits that cultural traits, innovations, and societal characteristics originate from a central point and then spread to other regions, making the statement true.
\end{enumerate}

\section*{Long Form Response}
\begin{enumerate}[label=\arabic*.]
\setcounter{enumi}{10}
\item \textbf{Answer:} The significance of the waterlogged soil at the Schöningen site in Germany lies in its ability to preserve ancient wooden spears that date back 300,000 years. The waterlogged conditions create an environment that prevents the decay of organic materials, allowing these wooden artifacts to survive for thousands of years. Without such unique soil conditions, these historical objects would likely have rotted away and been lost to time. This preservation offers valuable insights into the lives and hunting practices of early humans, making Schöningen an essential site for understanding human history and development.
\\ \textit{Note:} correct because it emphasizes how the waterlogged soil at the Schöningen site creates an anaerobic environment that inhibits the decay of organic materials, thereby allowing the ancient wooden spears to be preserved for study and understanding of early human life.
\item \textbf{Answer:} The Minoan civilization, which flourished on the island of Crete around 2000 BCE, emerged primarily as a local development influenced by the island's unique geography and resources. Although Crete was populated by various groups, including early Greeks and other Mediterranean peoples, it was the distinct adaptation of local conditions that fostered the growth of Minoan culture. The island's access to the sea facilitated trade and the exchange of ideas, while its agricultural potential supported a prosperous society. As a result, the Minoans developed advanced architecture, art, and governance systems that reflected their local innovations rather than solely external influences. Overall, the origins of Minoan civilization highlight the importance of Crete's geographical and cultural context in shaping a unique and influential culture in ancient history.
\\ \textit{Note:} correct because it emphasizes the significance of Crete's unique geography and resources in fostering the local development of Minoan civilization, while also acknowledging the role of diverse populations in shaping its cultural evolution through trade and agricultural practices.
\end{enumerate}

\vfill
\noindent\textit{End of Answer Key}
\end{document}